% sin solucion en decision -> optimo en relajacion
\begin{tikzpicture}[
    scale=0.85,
    transform shape,
    celda/.style={draw=gray!60, fill=gray!10, line width=0.6pt},
    verde/.style={circle, fill=green!60!black, draw=green!40!black, minimum size=0.65cm, inner sep=0pt},
    celeste/.style={circle, fill=cyan!75!blue!80, draw=cyan!50!black, minimum size=0.65cm, inner sep=0pt},
    tachado/.style={circle, draw=red!80!black, line width=1.3pt, densely dashed, minimum size=0.7cm, inner sep=0pt}
]

\newcommand{\cruz}[2]{
    \node[tachado] at (#1, #2) {};
    \draw[red!80!black, line width=1.4pt] (#1-0.21, #2-0.21) -- (#1+0.21, #2+0.21);
    \draw[red!80!black, line width=1.4pt] (#1-0.21, #2+0.21) -- (#1+0.21, #2-0.21);
}

% INSTANCIA INFACTIBLE EN DECISIÓN 
\begin{scope}[shift={(0,0)}]
    \draw[celda] (0,0) rectangle (3,4);
    \draw[celda] (0,0) grid (3,4);
    
    % Fila 1 (y=3.5): Verde en col 1, Celeste en col 2
    \node[verde] at (0.5, 3.5) {};
    \node[celeste] at (1.5, 3.5) {};
    % Fila 2 (y=2.5): Solo Celeste en col 1
    \node[celeste] at (0.5, 2.5) {};
    % Fila 3 (y=1.5): Solo Verde en col 2
    \node[verde] at (1.5, 1.5) {};
    % Fila 4 (y=0.5): Verde en col 3
    \node[verde] at (2.5, 0.5) {};
    
    \node[font=\footnotesize\sffamily] at (1.5, -0.45) {Tablero original};
\end{scope}

%  Seleccion de piedras removidas ÓPTIMA
\begin{scope}[shift={(4.2,0)}]
    \draw[celda] (0,0) rectangle (3,4);
    \draw[celda] (0,0) grid (3,4);
    
    % Fila 1 (queda vacía para salvar filas 2 y 3)
    \node[verde] at (0.5, 3.5) {}; \cruz{0.5}{3.5}
    \node[celeste] at (1.5, 3.5) {}; \cruz{1.5}{3.5}
    % Fila 2
    \node[celeste] at (0.5, 2.5) {};
    % Fila 3
    \node[verde] at (1.5, 1.5) {};
    % Fila 4
    \node[verde] at (2.5, 0.5) {};
    
    \node[font=\footnotesize\sffamily] at (1.5, -0.45) {Removidas óptimas};
\end{scope}

% SOLUCIÓN RELAJADA (3/4 FILAS
\begin{scope}[shift={(8.4,0)}]
    \draw[celda] (0,0) rectangle (3,4);
    \draw[celda] (0,0) grid (3,4);
    
    % Fila 1 vacía (y=3.5) - resaltada suavemente en rojo
    \draw[fill=red!8, draw=red!40, dashed] (0,3) rectangle (3,4);
    
    % Fila 2
    \node[celeste] at (0.5, 2.5) {};
    % Fila 3
    \node[verde] at (1.5, 1.5) {};
    % Fila 4
    \node[verde] at (2.5, 0.5) {};
    
    \node[font=\footnotesize\sffamily] at (1.5, -0.45) {Óptimo ($3$ de $4$ filas)};
\end{scope}

\end{tikzpicture}

